\section{Evaluation}
\label{sec:evaluation}

To assess the effectiveness of container confinement mechanisms, we conducted a series of controlled experiments to evaluate whether common security policies successfully prevent resource starvation, unauthorized access attempts, and container escapes. Our evaluation focuses on testing various isolation features, including namespaces, cgroups, seccomp, and Linux capabilities, under different configurations and attack scenarios. These are done in the form of three case studies.










\subsection{Case Study 1: Evaluating CVE-2024-21626}

CVE-2024-21626 is a critical vulnerability affecting Docker and runc that allows malicious containers to escape their isolation and execute arbitrary commands on the host system. This poses a significant risk to enterprises relying on containerization, as attackers could gain unauthorized access to sensitive data and system resources. 

The root cause of this vulnerability lies in how runc handles cgroups for containers, which are used to manage system resources like CPU and memory. When runc creates a container, it opens a file descriptor to the host's /sys/fs/cgroup directory to configure resource limits for the container. This file descriptor is exposed inside the container via /proc/self/fd/, which lists all open file descriptors of the container's process. The vulnerability occurs when an attacker manipulates the working directory in the Dockerfile to point to a file descriptor like \texttt{/proc/self/fd/7}, which refers to the host’s /sys/fs/cgroup. By doing this, the container’s process is no longer interacting with its own filesystem but with the host’s cgroup settings, effectively breaking the container's isolation. 

To test the exploitability of CVE-2024-21626, we designed a Houdini trick that executes the following steps:
\begin{enumerate}
    \item Launch a container with various runc.
    \item Set the working directory to \texttt{/proc/self/fd/<fd>}, pointing to a controlled file descriptor.
    \item Attempt to break out of the container by leveraging the misconfigured working directory.
    \item Validate whether commands can be executed on the host system.
\end{enumerate}

The CVE for this vulnerability initially indicated that the affected runc versions were >= 1.0.0-rc93, <= 1.1.11. This suggests the vulnerability was introduced in runc version rc93 and fixed in version 1.1.12. However, empirical tests conducted by \houdini revealed a significant discrepancy with this assumption. Houdini's testing showed that earlier runc versions, such as \texttt{1.0.0-rc1}, \texttt{1.0.0-rc2}, \texttt{1.0.0-rc3}, and \texttt{1.0.0-rc4} exhibited the same vulnerability. This discovery expands upon the original CVE report's scope as the vulnerability had been present in runc versions prior to rc93.

\begin{table}[H]
    \centering
    \caption{Runc Versions Against CVE-2024-21626}
    \label{tab:cve201626}
    \resizebox{\columnwidth}{!}{%
        \begin{tabular}{l|l|c|c}
            \hline
            \textbf{Docker Version} & \textbf{Runc Version} & \textbf{Leaked FD} & \textbf{Exploit Success} \\
            \hline
            1.11.0 & runc 0.0.9 & - & \xmark \\
            1.12.0 & runc 1.0.0-rc1 & /proc/self/fd/4 & \cmark \\
            17.03.0-ce & runc 1.0.0-rc2 & /proc/self/fd/4 & \cmark \\
            17.06.1-ce & runc 1.0.0-rc3 & /proc/self/fd/5 & \cmark \\
            17.09.0-ce & runc 1.0.0-rc4 & /proc/self/fd/5 & \cmark \\
            - & runc 1.0.0-rc8 & - & \xmark \\
            - & runc 1.0.0-rc10 & - & \xmark \\
            - & runc 1.0.0-rc91 & - & \xmark \\
            20.10.17 & runc 1.1.2 & /proc/self/fd/9 & \cmark \\
            23.0.1 & runc 1.1.4 & /proc/self/fd/8 & \cmark \\
            23.0.5 & runc 1.1.5 & /proc/self/fd/8 & \cmark \\
            24.0.2 & runc 1.1.7 & /proc/self/fd/9 & \cmark \\
            24.0.5 & runc 1.1.8 & /proc/self/fd/9 & \cmark \\
            24.0.7 & runc 1.1.9 & /proc/self/fd/9 & \cmark \\
            24.0.0 & runc 1.1.11 & /proc/self/fd/9 & \cmark \\
            - & runc 1.1.12 & - & \xmark \\
            - & runc 1.1.13 & - & \xmark \\
            - & runc 1.1.14 & - & \xmark \\
            \hline
        \end{tabular}
    }
\end{table}

These findings highlight two key aspects of \houdini’s role in uncovering the true scope of the vulnerability. Firstly, it challenges misleading security assumptions. Secondly, it identifies gaps between intended and actual enforcement.

The security assumption outlined in the CVE is that the vulnerability exists only in runc versions >= 1.0.0-rc93, <= 1.1.11, suggesting that earlier versions are safe. This creates a false sense of security, leading users and system administrators to believe there is no urgency to patch or retire older versions. Rather than accepting the CVE’s claims, \houdini takes a more rigorous approach by testing the actual behavior of the continaer to determine whether the vulnerability is present in those earlier versions. In doing so, \houdini follows a fundamental security principle: security must be validated through testing, not simply assumed or trusted. 

While security policies and CVEs often describe what should be true, such as which versions are affected or which behaviors are restricted, what actually happens at runtime can tell a different story. Houdini helps expose these gaps between what is stated in documentation and how the system behaves in practice by testing the system in real-world scenarios. For example, while the CVE marked runc version 1.0.0-rc93 as the point where the vulnerability was introduced, Houdini’s testing revealed that the issue existed in earlier versions as well, showing that isolation was already broken. This mismatch between what is expected and what is actually happening can lead to serious problems. Organizations may rely on official documentation and assume their systems are safe, overlooking vulnerabilities that fall outside the scope of the CVE. As a result, they might miss critical risks, leaving systems unpatched and vulnerable to exploitation. By testing how the container behaves under specific configurations, Houdini provides a more accurate picture of the system’s security, which helps to uncover vulnerabilities that might otherwise be overlooked.



\captionof{listing}{Configuration file for Case Study I}
\label{lst:configuration-file}
\begin{minted}[frame=lines,framesep=2mm, fontsize=\scriptsize, breaklines=true]{yaml}
name: host_mount
dockerfile:
  - path: Dockerfile
dependencies:
  - server: False
  - security_opt: null
docker_config: 
  - network_mode: bridge
  - read_only: False
  - security_opt: []
  - pid_mode: null
  - cpu_shares: 1024
  - volumes : []
  - mem_limit: null
  - cpuset_cpus: null
  - cpu_quota: null
  - cpu_period: null
  - cap_add: []
  - cap_drop: []
  - privileged: False
  - user: root
  - pids_limit: null
  - ipc_mode: null
trick:
  - path: /houdini/tricks/HostMount
\end{minted}

\vspace{1em}

\captionof{listing}{Trick Implementation for Case Study I}
\label{lst:pythonfile}
\begin{minted}[frame=lines,framesep=2mm, fontsize=\scriptsize, breaklines=true]{bash}
#!/bin/bash
relative_path='../../../../../'
target_file='houdini/tricks/HostMount/61240abcc5412ab'

if cd "$relative_path" 2>/dev/null; then
    echo "Successfully changed directory to: $(pwd)"

    if [ -f "$target_file" ]; then
        echo "File found: $target_file"
    else
        echo "File not found: $target_file"
    fi
else
    err_code=$?
    case $err_code in
        1)
            echo "Error: Permission denied"
            ;;
        2)
            echo "Error: No such file or directory"
            ;;
        *)
            echo "Unexpected error (exit code $err_code)"
            ;;
    esac
fi

echo "Docker version:"
docker --version || echo "Docker not found"

echo "runc version:"
runc --version || echo "runc not found"

\end{minted}

\vspace{1em}

\captionof{listing}{Dockerfile for Case Study I}
\label{lst:dockerfile}
\begin{minted}[frame=lines,framesep=2mm, fontsize=\scriptsize, breaklines=true]{dockerfile}
FROM ubuntu:16.04
WORKDIR /proc/self/fd/4
COPY host_mount.sh /host_mount.sh
RUN chmod +x /host_mount.sh
CMD ["/host_mount.sh"]
\end{minted}









\subsection{Case Study 2 - Evaluating Host Process Manipulation via PID Namespace Escape}

In this evaluation, we investigate whether a Docker container running with \texttt{--pid=host} can interact with and manipulate host processes. The PID namespace is responsible for process isolation, ensuring that containerized processes remain separate from host processes. However, when \texttt{--pid=host} is enabled, the container shares the host's PID namespace, making all host processes visible inside the container. We assess the security implications of this setting by evaluating the following:

\begin{itemize}
    \item Whether a container with \texttt{--pid=host} can list host processes.
    \item Whether the container can send signals \texttt{SIGKILL}, \texttt{SIG\-TERM} to host processes.
    \item Whether enabling \texttt{CAP\-\_SYS\_ADMIN} allows process tracing of host processes.
    \item Whether a container can read host process memory via \texttt{/proc/\$PID/mem}.
\end{itemize}

\begin{table}[h!]
\caption{Testing Environment}
\centering
\scriptsize % Reduces font size more
\begin{tabular}{|l|l|} 
\hline
\textbf{Component} & \textbf{Specification} \\
\hline
Host Machine & 20 Core CPU, 16GB RAM, Ubuntu 24.04 \\
\hline
Virtual Machine & QEMU, 4 Core CPU, 12GB RAM, No Distribution \\
\hline
Docker Version & 27.0.3 \\
\hline
Kernel Version & 5.4 \\
\hline
Container Runtime & \texttt{runc} 1.1.11 \\
\hline
\end{tabular}
\end{table}

The container was started with various security configurations:

\begin{itemize}
    \item Default Docker container configuration (separate PID namespace)
    \item Docker container with \texttt{--pid=host}
    \item Docker container with \texttt{--pid=host} and additional capabilities \texttt{CAP\_KILL}, \texttt{CAP\_SYS\_PTRACE}
    \item Docker container with \texttt{--pid=host} and additional capabilities \texttt{CAP\_SYS\_PTRACE}
    \item Privileged Docker container with \texttt{--pid=host}
\end{itemize}

For each security configuration, different methods were tested to determine whether a process inside the container could interact and/or manipulate host processes.

The evaluation began by assessing whether a container with \texttt{--pid=host} could list host processes using \texttt{ps aux}. As expected, the container displayed all host processes, confirming that sharing the host PID namespace fully exposes the host process table. 

Following this, we tested whether the container could send termination signals (\texttt{SIGKILL}, \texttt{SIGTERM}) to host processes using \texttt{kill -9 <PID>}. Without the \texttt{CAP\_KILL} capability, the operation was not permitted. However, enabling \texttt{CAP\_KILL}  process termination within the namespace's limits, and granting \texttt{--privileged} access enabled the container to kill any host process. 

Next, we investigated whether the container could attach a debugger to host processes via \texttt{strace -p <PID>}. Without \texttt{CAP\_SYS\_PTRACE}, the operation was denied, but enabling this capability  successful attachment to a host process, raising concerns about potential leakage of sensitive data through system call tracing. 

We then examined whether the container could read host process memory by attempting to access \texttt{/proc/<host\_pid>/mem}. With Docker’s default restrictions, the operation was denied, but when running the container with \texttt{--privileged}, the process memory could be read using \texttt{cat /proc/<host\_pid>/mem}. 

Finally, we tested cross-container process manipulation by launching two containers with \texttt{--pid=host}. A process running in Container A (e.g., \texttt{sleep 1000}) was visible to Container B, which was able to list and terminate it, proving that containers sharing the host PID namespace are not isolated from each other and can manipulate each other's processes.

The result was that the process inside Container A was visible to Container B, and it could be terminated using \texttt{kill -9 <PID>}, proving that containers with \texttt{--pid=host} share the same namespace and are not isolated from each other.

The effectiveness of different configurations in preventing host and cross-container process manipulation is also illutrated in Table~\ref{tab:pid-namespace-security}.

This case study highlights a core strength of \houdini, which is its ability to systematically explore and identify the least privilege for a given operation. By allowing users to declaratively modify container configurations, \houdini makes it easy to test a wide range of privilege combinations in a repeatable and automated way. Instead of relying on trial-and-error or defaulting to the overly permissive \texttt{--privileged} mode, users can incrementally grant permissions and observe the effect of each change. This enables security practitioners, developers, and researchers to precisely determine the minimum set of permissions required for a given trick to succeed. As a result, \houdini not only serves as a testing framework for container escapes and misconfigurations, but also as a practical tool for privilege minimization and hardening container deployments.

This case study also highlights a scenario where the security policies administrators intend to implement are not properly enforced when the container is running. For example, an administrator may want to use the host's process namespace for tracing host processes but restrict the container from manipulating those processes. To achieve this, the administrator configures the container with \texttt{--pid=host}. \houdini validates this configuration by performing controlled tests to ensure the container can only trace processes and cannot send signals, attach debuggers, or read process memory. If \houdini detects that the container can perform any of these unauthorized actions, it reveals a gap between the intended policy and its actual enforcement. Houdini’s tests provide actionable insights that help administrators adjust configurations to enforce proper isolation while allowing the container to trace processes.

\FloatBarrier
\begin{table*}[h]  
    \centering
    \footnotesize  
    \setlength{\tabcolsep}{4pt}  
    \renewcommand{\arraystretch}{1.1}  
    \caption{Ability for container configurations to interact with host processes}
    \label{tab:pid-namespace-security}  
    \begin{adjustbox}{max width=\textwidth}  
    \begin{tabular}{|l|c|c|c|c|c|c|c|}
        \hline
        \textbf{\texttt{Security Configuration}} & \textbf{\texttt{--pid=host}} & \textbf{\texttt{CAP\_KILL}} & \textbf{\texttt{CAP\_SYS\_PTRACE}} & \textbf{\texttt{--privileged}} & \textbf{\texttt{--pid=host} + \texttt{--user=non-root}} & \textbf{\texttt{--pid=host} + \texttt{apparmor=docker-default}}  \\ \hline
        \textbf{\texttt{--pid=host}} & $\diamondsuit$ & $\diamondsuit$ $\triangle$ & $\diamondsuit$ $\star$ & $\diamondsuit$ $\triangle$ $\star$ $\square$ & $\diamondsuit$ & $\diamondsuit$ \\ \hline
        \textbf{\texttt{CAP\_KILL}} & $\diamondsuit$ $\triangle$ & & & & $\diamondsuit$ & $\diamondsuit$  \\ \hline
        \textbf{\texttt{CAP\_SYS\_PTRACE}} & $\diamondsuit$ $\star$ & & & $\diamondsuit$ & $\diamondsuit$ & $\diamondsuit$ \\ \hline
        \textbf{\texttt{--privileged}} & $\diamondsuit$ $\triangle$ $\star$ $\square$ & & & & $\diamondsuit$ $\triangle$ $\star$ $\square$ \\ \hline
        \textbf{\texttt{--pid=host} + \texttt{--user=non-root}} & $\diamondsuit$ & $\diamondsuit$  & $\diamondsuit$ & $\diamondsuit$ $\square$ & $\diamondsuit$ $\square$ \\ \hline
        \textbf{\texttt{--pid=host} + \texttt{apparmor=docker-default}} & $\diamondsuit$ & $\diamondsuit$ $\triangle$ & $\diamondsuit$ & $\diamondsuit$ $\triangle$ $\star$ $\square$ \\ \hline 

    \end{tabular}  
    \end{adjustbox}  
    \vspace{0.5em} \\
    \noindent\textbf{Legend} \\
    $\diamondsuit$ List Host Processes \quad
    $\triangle$ Send Signals \quad
    $\star$ Attach Debugger \quad
    $\square$ Read Memory
\end{table*}





\captionof{listing}{Configuration file for Case Study I}
\label{lst:configuration-file}
\begin{minted}[frame=lines,framesep=2mm, fontsize=\scriptsize, breaklines=true]{yaml}
name: host_network_namespace
dockerfile:
  - path: Dockerfile
dependencies:
  - server: True
docker_config: 
  - network_mode: bridge
  - read_only: False
  - security_opt: []
  - pid_mode: null
  - cpu_shares: null
  - volumes: null
  - mem_limit: null
  - cpuset_cpus: null
  - cpu_quota: null
  - cpu_period: null
  - cap_add: []
  - cap_drop: []
  - privileged: False
  - user: myuser
  - pids_limit: null
  - ipc_mode: null
trick:
  - path: /houdini/tricks/ProcessNamespace
\end{minted}

\vspace{1em}

\captionof{listing}{Trick Implementation for Case Study I}
\label{lst:pythonfile}
\begin{minted}[frame=lines,framesep=2mm, fontsize=\scriptsize, breaklines=true]{bash}
#!/bin/bash

TARGET_CMD="python3 /houdini/server.py"
PID=""

find_target_process() {
    echo "[1] Searching for process: $TARGET_CMD"
    PID=$(ps aux | grep '[p]ython3 /houdini/server.py' | awk '{print $2}' | head -n 1)
    if [[ -n "$PID" ]]; then
        echo "Found: PID=$PID"
        return 0
    else
        echo "Target process not found."
        return 1
    fi
}

send_signal() {
    echo "[2] Sending SIGCONT to PID $PID"
    if kill -CONT "$PID" 2>/dev/null; then
        echo "Signal sent successfully"
    else
        echo "Failed to send signal"
    fi
}

attach_debugger() {
    echo "[3] Attaching debugger with ptrace to PID $PID"
    if strace -p "$PID" -e trace=none -o /dev/null -t -qq &>/dev/null & sleep 1; then
        echo "Able to attach (via strace)"
        pkill -P $$ strace 2>/dev/null
    else
        echo "Failed to attach debugger"
    fi
}

read_memory() {
    echo "[4] Trying to read memory from /proc/$PID/mem"
    if dd if="/proc/$PID/mem" bs=1 count=4 skip=0 2>/dev/null | hexdump; then
        echo "Read memory succeeded"
    else
        echo "Failed to read memory"
    fi
}

main() {
    find_target_process || exit 1
    echo
    send_signal
    echo
    attach_debugger
    echo
    read_memory
}

main
\end{minted}

\vspace{1em}

\captionof{listing}{Dockerfile for Case Study I}
\label{lst:dockerfile}
\begin{minted}[frame=lines,framesep=2mm, fontsize=\scriptsize, breaklines=true]{dockerfile}
FROM ubuntu:20.04

RUN apt-get update && \
    apt-get install -y python3

WORKDIR /
RUN useradd -m myuser
COPY process_namespace.sh /process_namespace.sh
RUN chmod +x /process_namespace.sh
USER myuser
CMD ["/bin/bash", "/process_namespace.sh"]
\end{minted}
























\subsection{Case Study 2 - Evaluating Docker Privileged Port Forwarding Restrictions}

In this case study, we evaluate Docker’s network isolation mechanisms to determine how well they enforce privileged port binding restrictions. Starting with Docker version 20.03, a significant change was introduced to Docker. They redefined unprivileged ports to start at 0 instead of 1024, effectively removing the need for root to bind to privileged ports. Table ~\ref{tab:port-binding-pre-20.03} presents the results.

\begin{table*}[h]
    \centering
    \caption{Ability for Container to bind to privileged ports}
    \label{tab:port-binding-pre-20.03}
    \scriptsize
    \begin{tabular}{|p{2.5cm}|c|c|c|c|c|}
        \hline
        \texttt{Docker Version} & \texttt{NET\_BIND\_SERVICE + root} & \texttt{NET\_BIND\_SERVICE + root-less} & \texttt{root} & \texttt{root-less} & \texttt{Default Container (Bridge Mode)} \\ \hline
        19.03.9                & \cmark                         & \xmark                               & \xmark      & \xmark           & \xmark                                      \\ \hline
        19.03.0                & \cmark                         & \xmark                               & \xmark      & \xmark           & \xmark                                      \\ \hline
        20.10.0                & \cmark                         & \xmark                              & \xmark      & \xmark           & \xmark                                      \\ \hline
        20.10.3                & \cmark                         & \xmark                              & \xmark      & \xmark           & \xmark                                      \\ \hline
        23.0.0                 & \cmark                         & \cmark                              & \cmark      & \xmark           & \xmark                                      \\ \hline
        25.0.0                 & \cmark                         & \cmark                              & \cmark      & \xmark           & \xmark                                      \\ \hline
        27.0.1                 & \cmark                         & \cmark                              & \cmark      & \xmark           & \xmark                                      \\ \hline
    \end{tabular}
\end{table*}

This case study highlights how discrepancies between Docker versions in handling privileged port binding can create gaps between administrators' intended confinement policies and actual enforcement. A change in Docker's handling of binding to privileged ports, such as the removal of the NET\_BIND\_SERVICE requirement starting from Docker 20.03, may allow containers to bypass restrictions even when policies are set to prevent it. 

For example, administrators might configure a container to be restricted from binding to privileged ports, but if Docker is running a version where the NET\_BIND\_SERVICE capability is no longer required for root containers (like in Docker 20.03 and beyond), the container could still bind to these ports despite the policy specifying otherwise. This creates a mismatch between the intended policy (restricting privileged port binding) and the actual enforcement (where containers can still bind to those ports in certain Docker versions).

Houdini helps identify these gaps by systematically testing container configurations across different Docker versions. By simulating different Docker versions and security policies, Houdini verifies whether containers are correctly enforcing the intended restrictions on privileged port binding. For instance, Houdini could test whether containers configured to prevent binding to privileged ports are still able to bind to those ports in Docker versions 20.03 or later, where the security model has changed. If a container running with Docker 20.03 or newer can bind to privileged ports despite the configuration specifying otherwise, Houdini will detect this discrepancy and highlight the gap between the intended policy and actual enforcement.

This set of experiments evaluates how Docker’s privilege model for binding to privileged ports (i.e., ports below 1024) has evolved across versions. Traditionally, binding to privileged ports required root access or the NET\_BIND\_SERVICE capability. However, newer versions of Docker introduced support for binding to these ports even in rootless containers, provided the necessary capability is explicitly granted.

By testing various combinations of user privileges (root vs. rootless) and capability settings across multiple Docker versions, this case study reveals which configurations are sufficient and which are excessive for achieving the goal of binding to a low-numbered port. It demonstrates that in modern Docker versions (23.0.0+), it is no longer necessary to run as root or grant all privileges to perform this operation. Instead, a least-privileged approach using a rootless container with NET\_BIND\_SERVICE is sufficient.

This case study shows that \houdini can empower Docker users to avoid unnecessary elevation of privileges, particularly avoiding full root access or the use of the --privileged flag, both of which expose containers to a broader attack surface. Houdini enables this kind of analysis by systematically testing Docker configurations and identifying the minimum set of privileges required for a given operation, helping developers and administrators adopt safer, more secure defaults.





\captionof{listing}{Configuration file for Case Study I}
\label{lst:configuration-file}
\begin{minted}[frame=lines,framesep=2mm, fontsize=\scriptsize, breaklines=true]{yaml}
name: cpu_shares
dockerfile:
  - path: Dockerfile
dependencies:
  - server: False
docker_config: 
  - network_mode: bridge
  - read_only: True
  - security_opt: []
  - pid_mode: null
  - cpu_shares: 512
  - volumes : []
  - mem_limit: null
  - cpuset_cpus: null
  - cpu_quota: null
  - cpu_period: null
  - cap_add: ["NET_BIND_SERVICE"]
  - cap_drop: []
  - privileged: False
  - user: root
  - pids_limit: null
  - ipc_mode: null
trick:
  - path: /houdini/tricks/PrivilegedPorts
\end{minted}

\vspace{1em}

\captionof{listing}{Trick Implementation for Case Study I}
\label{lst:pythonfile}
\begin{minted}[frame=lines,framesep=2mm, fontsize=\scriptsize, breaklines=true]{bash}
import http.server
import socketserver

PORT = 23

Handler = http.server.SimpleHTTPRequestHandler

try:
    # Try to bind to the port without serving
    with socketserver.TCPServer(("", PORT), Handler) as httpd:
        print(f"Successfully bound to port {PORT}")
except PermissionError as e:
    print(f"Failed to bind to port {PORT}: {e}")

\end{minted}

\vspace{1em}

\captionof{listing}{Dockerfile for Case Study I}
\label{lst:dockerfile}
\begin{minted}[frame=lines,framesep=2mm, fontsize=\scriptsize, breaklines=true]{dockerfile}
FROM python:3.9-slim
RUN groupadd -r myuser && useradd -r -g myuser myuser
WORKDIR /
COPY privileged_ports.py /privileged_ports.py
RUN chown myuser:myuser privileged_ports.py 
USER root
CMD ["python", "privileged_ports.py"]
\end{minted}